%%%%%%%%%%%%%%%%%
% CV tailored for Thales - IS Data Engineer / Data Governance / Data Catalog / Datalake
%%%%%%%%%%%%%%%%%

\documentclass[8pt,a4paper,ragged2e,withhyper]{../_shared/altacv}

\geometry{left=1.25cm,right=1.25cm,top=.5cm,bottom=1.cm,columnsep=1.2cm}

\usepackage{paracol}
\usepackage{fontawesome5}
\usepackage{hyperref}

\iftutex 
  \setmainfont{Roboto Slab}
  \setsansfont{Lato}
  \renewcommand{\familydefault}{\sfdefault}
\else
  \usepackage[rm]{roboto}
  \usepackage[defaultsans]{lato}
  \renewcommand{\familydefault}{\sfdefault}
\fi

\definecolor{SlateGrey}{HTML}{2E2E2E}
\definecolor{LightGrey}{HTML}{666666}
\definecolor{DarkPastelRed}{HTML}{8AC0D9}
\definecolor{PastelRed}{HTML}{00B0DA}
\definecolor{GoldenEarth}{HTML}{B4AD0C}

\colorlet{name}{black}
\colorlet{tagline}{PastelRed}
\colorlet{heading}{DarkPastelRed}
\colorlet{headingrule}{GoldenEarth}
\colorlet{subheading}{PastelRed}
\colorlet{accent}{PastelRed}
\colorlet{emphasis}{SlateGrey}
\colorlet{body}{LightGrey}

\definecolor{violet}{HTML}{5A2582}
\definecolor{rouge}{HTML}{F53B3B}
\definecolor{noir}{HTML}{00232C}
\definecolor{bleu}{HTML}{00B0DA}
\definecolor{rose}{HTML}{F831A0}
\definecolor{jaune}{HTML}{FFEC00}
\definecolor{vert}{HTML}{34AD6C}

\renewcommand{\namefont}{\Huge\rmfamily\bfseries}
\renewcommand{\personalinfofont}{\footnotesize}
\renewcommand{\cvsectionfont}{\LARGE\rmfamily\bfseries}
\renewcommand{\cvsubsectionfont}{\large\bfseries}
\renewcommand{\cvItemMarker}{{\small\textbullet}}
\renewcommand{\cvRatingMarker}{\faCircle}

\begin{document}

\name{BOILEAU Guillaume}
\tagline{IS Data Engineer | Data Governance, Data Quality, Data Cataloging, Datalake \& Application Support}

\photoL{2.8cm}{../_shared/moi}

\personalinfo{%
  \email{guillaume.boileaupro@gmail.com}
  \phone{+33 6 59 94 85 84}
  \location{Alpes-Maritimes, FRANCE}
  \linkedin{guillaume-boileau-phd-891b81265}
  \github{guillaume-boileaupro}
}

\makecvheader
\vspace{-0.3cm}

\AtBeginEnvironment{itemize}{\small}
\columnratio{0.64}

\begin{paracol}{2}

\cvsection{Experience}

\cvevent{\textbf{Software Engineer - Operational Data, Application Support \& Industrial Interfaces}}{VEV Platform Services France}{2025 -- 2026}{Valbonne, France}

\textbf{Context:} Development and support of software services for an EV charging infrastructure CPMS platform connected to charging stations, operational data flows, industrial interfaces and field equipment.

\textbf{Objective:} Support reliable operational data exchanges, application behaviour analysis, production diagnostics and coordination between software services, infrastructure constraints and user needs.

\begin{itemize}
  \item \textbf{Operational data flows:} Implemented and monitored FTP-based data exchange services used for infrastructure operation, technical follow-up and platform data handling.
  \item \textbf{Application support:} Investigated production issues involving data exchanges, service behaviour, interface continuity, operational reliability and platform constraints.
  \item \textbf{Data consistency:} Checked consistency between CPMS platform services, operational data flows, charging station behaviour and expected business logic.
  \item \textbf{Industrial interfaces:} Worked on services connected to EV charging stations, electrical equipment status, field data and maintenance-related diagnostics.
  \item \textbf{Incident analysis:} Used logs, application behaviour analysis and technical diagnostics to identify root causes and support service continuity.
  \item \textbf{User support:} Clarified issues with technical and operational stakeholders, structured feedback and documented observed behaviours.
  \item \textbf{Agile environment:} Used Zenhub for task tracking, backlog follow-up and coordination within an iterative software development environment.
  \item \textbf{Software development:} Developed web and backend services using TypeScript, Angular and Node.js in a production-oriented industrial context.
\end{itemize}

\divider

\cvevent{\textbf{Senior R\&D Engineer - Data Integration, Data Quality \& Validation Platform}}{Laboratoire Lagrange, Observatoire de la Côte d'Azur, CNES}{2023 -- 2025}{Nice, France}

\textbf{Context:} Engineering and R\&D activities for the LISA ESA/NASA space mission, involving large-scale simulations, heterogeneous model outputs, technical validation, data structuring and international coordination.

\textbf{Objective:} Structure, validate, document and exploit complex datasets while supporting contributors, comparison campaigns, reproducibility, technical decision-making and shared data practices.

\begin{itemize}
  \item \textbf{Data engineering:} Developed Python-based workflows to ingest, structure, process, compare and validate heterogeneous simulation outputs from multiple contributors.
  \item \textbf{Data quality:} Implemented checks focused on data consistency, completeness, traceability, reproducibility, usability and comparison of technical outputs.
  \item \textbf{Data cataloging logic:} Structured model configurations, simulation products, metadata-like information, analysis outputs and validation evidence for reuse and review.
  \item \textbf{Validation platform:} Built CATpyLISA, a Python platform for model integration, data comparison, technical validation and review-oriented analysis workflows.
  \item \textbf{Governance and good practices:} Helped define shared practices for data usage, result comparison, documentation, reproducibility and technical review.
  \item \textbf{Stakeholder interface:} Acted as a technical interface between users, contributors and project stakeholders to clarify needs, constraints and expected outputs.
  \item \textbf{Technical documentation:} Used Confluence to organise technical documentation, coordination notes, validation follow-up and knowledge sharing.
  \item \textbf{Roadmap support:} Identified user needs, workflow limitations, technical improvements and data issues to support the evolution of internal tools.
  \item \textbf{Technical synthesis:} Produced structured reports, validation summaries, review material and decision-support documents for complex data workflows.
  \textcolor{DarkPastelRed}{\href{https://gitlab.oca.eu/gboileau/lisa_l3_tools}{\bf GitLab:CATpyLISA}}
\end{itemize}

\divider

\cvevent{\textbf{Data Scientist - Statistical Pipelines, Data Quality \& Technical Analysis}}{Universiteit Antwerpen, Belgium}{2021 -- 2023}{Antwerp, Belgium}

\textbf{Context:} Data analysis and performance studies for precision instruments in a high-requirement technical environment linked to gravitational-wave detector performance.

\textbf{Objective:} Build robust data analysis pipelines, validate complex datasets and produce documented results to support understanding of system performance and technical risks.

\begin{itemize}
  \item \textbf{Data pipelines:} Developed Python-based workflows for noise source identification, characterization, mitigation and sensitivity studies.
  \item \textbf{Data validation:} Processed heterogeneous and noisy datasets with attention to consistency, reproducibility, uncertainty and robustness of conclusions.
  \item \textbf{Statistical analysis:} Applied Bayesian inference and MCMC methods to extract reliable information from complex and incomplete data environments.
  \item \textbf{Technical reporting:} Transformed complex datasets and statistical results into structured conclusions for technical stakeholders.
  \item \textbf{Decision support:} Helped identify environmental and instrumental constraints affecting next-generation detector design and performance.
  \item \textbf{Technical communication:} Worked in English within international collaborations and multidisciplinary technical environments.
\end{itemize}

\divider

\cvevent{\textbf{Junior R\&D Engineer - Simulation, Data Processing \& Validation}}{ARTEMIS, Observatoire de la Côte d'Azur}{2018 -- 2021}{Nice, France}

\textbf{Context:} R\&D activities for gravitational-wave mission preparation, involving signal analysis, simulation, technical documentation and characterization of complex low-SNR data environments.

\textbf{Objective:} Develop analysis methods, process complex simulation outputs, document results and support validation activities in demanding system environments.

\begin{itemize}
  \item \textbf{Simulation workflows:} Built simulation approaches for complex signal environments, uncertainty studies and large-scale datasets.
  \item \textbf{Data processing:} Processed, compared and interpreted complex outputs produced by numerical models and signal-processing pipelines.
  \item \textbf{Validation logic:} Developed methods for separating overlapping low-SNR signals and assessing data quality in challenging conditions.
  \item \textbf{Technical investigations:} Investigated noise sources through cross-sensor correlation analyses in diagnostic contexts.
  \item \textbf{Documentation:} Documented methods, assumptions, outputs and conclusions for collaborative review and reproducibility.
\end{itemize}

\switchcolumn

\cvsection{Profile for Thales}

\textbf{IS Data Engineer-oriented profile with strong experience in data engineering, data quality, data validation, technical documentation, application support and coordination of complex data workflows}. Background developed through major international collaborations including \textbf{LISA ESA/NASA}, \textbf{LIGO--Virgo--KAGRA} and \textbf{Einstein Telescope}, complemented by recent industrial software experience involving operational data flows, production diagnostics and connected equipment.

For the \textbf{IS Data Engineer / Data Galaxy / Datalake} role, I bring a strong ability to structure information, support data cataloging practices, improve data quality, contribute to governance logic, analyse incidents, support users and connect technical constraints with business needs.

\begin{itemize}
  \item Strong experience in data structuring, validation, consistency checks, reproducibility and traceability.
  \item Practical experience with complex data workflows involving heterogeneous sources, simulation products, operational data and technical outputs.
  \item Knowledge of AWS S3 principles, object storage, bucket-based organisation, data ingestion and datalake logic.
  \item Ability to support data cataloging and metadata-oriented practices through structured documentation, repositories and reusable workflows.
  \item Experience investigating application and production issues through logs, monitoring, technical diagnostics and root-cause analysis.
  \item Ability to act as an interface between users, software teams, technical contributors and operational stakeholders.
  \item Practical use of collaborative tools including Confluence for documentation, Zenhub for activity follow-up, Jira for issue tracking and OpenSearch/log-based analysis in operational contexts.
  \item Understanding of data governance principles: ownership, consistency, quality, completeness, traceability, documentation and controlled reuse.
  \item Strong analytical mindset for transforming complex technical information into clear summaries and decision-support material.
  \item Comfortable working with cloud-oriented data environments, object storage logic and datalake principles.
  \item Strong communication skills in English within international and multidisciplinary environments.
\end{itemize}

\cvsection{Languages}

\cvskill{English}{4}
\divider
\cvskill{Spanish}{2}
\divider

\medskip

\cvsection{Education}

\cvevent{PhD in Physics}{Université Côte d'Azur}{2018 -- 2021}{Nice, France}

\divider

\cvevent{Selective Graduate Program in Fundamental Physics}{Université Paris-Saclay}{2016 -- 2018}{Orsay, France}

\divider

\cvevent{MSc in Astronomy and Astrophysics}{Observatoire de Paris / Université Paris-Saclay}{2017 -- 2018}{Paris, France}

\divider

\cvevent{BSc in Fundamental Physics}{Université Paris-Saclay}{2015 -- 2016}{Orsay, France}

\cvsection{Professional Training}

\cvevent{Supervised Machine Learning - Regression \& Classification}{DeepLearning.AI - Andrew Ng}{2020}{Coursera}

\end{paracol}

\cvsection{Projects}

\cvevent{CATpyLISA - Data Integration, Cataloging Logic \& Validation Platform}{Python Platform for the LISA ESA/NASA Space Mission}{2023 -- 2025}{}

Development of a Python platform dedicated to large-scale model integration, data structuring, comparison, validation and technical review for LISA mission-level analysis workflows.

\textbf{Role:} Data engineering, technical coordination, data structuring, documentation support, validation reporting, consistency checks and interaction with multiple contributors.

\textbf{Relevance for IS Data Engineering:} Work involving heterogeneous datasets, structured repositories, traceability of assumptions and outputs, quality checks, reusable workflows, documentation, cataloging-like logic and knowledge capitalisation.

\divider

\cvevent{Operational Data Flows - EV Charging Platform}{Application Interfaces, Production Diagnostics and Data Exchange Services}{2025 -- 2026}{}

Contribution to operational data exchange mechanisms and software services within an EV charging infrastructure platform connected to field equipment and production constraints.

\textbf{Role:} Development and support of data exchange services, incident investigation, application behaviour analysis, operational feedback and coordination with technical stakeholders.

\textbf{Relevance for Application Support:} Work involving data consistency, interface continuity, logs, operational reliability, production diagnostics, user support and service follow-up.

\cvsection{Strengths}

\vspace{-.3cm}

\cvsubsection{Data Engineering, Governance \& Cataloging}

{\small
\cvtag{Data Engineering}
\cvtag{Data Governance}
\cvtag{Data Quality}
\cvtag{Data Cataloging}
\cvtag{Data Galaxy}
\cvtag{Datalake}
\cvtag{AWS S3}
\cvtag{Object Storage}
\cvtag{Bucket Organisation}
\cvtag{Data Ingestion}
\cvtag{Data Structuring}
\cvtag{Data Repositories}
\cvtag{Metadata}
\cvtag{Data Lineage}
\cvtag{Traceability}
\cvtag{Data Completeness}
\cvtag{Data Consistency}
\cvtag{Controlled Reuse}
}

\divider\smallskip
\vspace{-.3cm}

\cvsubsection{Application Support, Incident Analysis \& User Interface}

{\small
\cvtag{Application Support}
\cvtag{Incident Analysis}
\cvtag{Root-Cause Analysis}
\cvtag{Log Analysis}
\cvtag{Monitoring}
\cvtag{OpenSearch}
\cvtag{Production Diagnostics}
\cvtag{Operational Data Flows}
\cvtag{Interface Continuity}
\cvtag{User Support}
\cvtag{Business Needs}
\cvtag{Technical Support}
\cvtag{Service Reliability}
\cvtag{Issue Follow-Up}
\cvtag{Roadmap Support}
\cvtag{Best Practices}
}

\divider\smallskip
\vspace{-.3cm}

\cvsubsection{Software, Data Pipelines \& Technical Tools}

{\small
\cvtag{Python}
\cvtag{SQL}
\cvtag{Bash}
\cvtag{Linux}
\cvtag{Git}
\cvtag{GitLab}
\cvtag{GitHub}
\cvtag{Docker}
\cvtag{CI/CD}
\cvtag{ETL Concepts}
\cvtag{Data Pipelines}
\cvtag{REST API}
\cvtag{FTP}
\cvtag{TypeScript}
\cvtag{Angular}
\cvtag{Node.js}
\cvtag{JavaScript}
\cvtag{C}
\cvtag{Matlab}
}

\divider\smallskip
\vspace{-.3cm}

\cvsubsection{Coordination, Documentation \& Technical Analysis}

{\small
\cvtag{Technical Documentation}
\cvtag{Confluence}
\cvtag{Jira}
\cvtag{Zenhub}
\cvtag{Stakeholder Coordination}
\cvtag{Technical Coordination}
\cvtag{Reporting}
\cvtag{Synthesis Reports}
\cvtag{Decision Support}
\cvtag{Review Preparation}
\cvtag{Knowledge Capitalisation}
\cvtag{Process Improvement}
\cvtag{Validation Follow-Up}
\cvtag{Verification}
\cvtag{Reproducibility}
\cvtag{Complex Systems}
\cvtag{Space Systems}
}

\end{document}